\section*{Chapitre \ref{ch:b2:eruptive-styles} --- Éruptions volcaniques}

\begin{solution}{exo:b2:eruptive-styles:1}
(a) $\qty{1.3}{km^3} = \qty{1.3e9}{m^3}$, entre $10^9$ et $10^{10}$~: VEI~5.
(b) $\qty{2.9e8}{m^3}$~: VEI~4. (c) $\qty{2e6}{m^3}$, entre $10^6$ et $10^7$~:
VEI~2. (d) Le VEI~6 commence à \qty{e10}{m^3}, donc le volume apparent de \omterm{def:b2:eruptive-styles:tephra}{téphra} doit être d'au moins
\qty{10}{km^3}~; «~plus de \qty{5}{km^3}~» est une borne inférieure, non une contradiction.
\end{solution}

\begin{solution}{exo:b2:eruptive-styles:2}
(a) Cendre (plus fine que \qty{2}{mm}). (b) Un lapillus (\qtyrange{2}{64}{mm}).
(c) Un bloc~: plus grand que \qty{64}{mm}, anguleux, solide au moment de son éjection. (d) Une
\omterm{def:b2:eruptive-styles:tephra}{bombe volcanique}, une bombe fusiforme~: plus grande que \qty{64}{mm} et façonnée à l'état fondu
en vol.
\end{solution}

\begin{solution}{exo:b2:eruptive-styles:3}
(a) $0.304\times(10^6)^{0.241} = 0.304\times27.9 = \qty{8.5}{km}$;
$0.304\times(10^8)^{0.241} = 0.304\times84.7 = \qty{25.8}{km}$.
(b) $Q = (35/0.304)^{1/0.241} = 115.1^{4.149} = \qty{3.6e8}{kg/s}$, et $I =
8.55 + 3 = 11.6$. (c) $Q \propto H^{4.15}$, donc une hauteur surestimée de \qty{20}{\percent} surestime
$Q$ d'un facteur $1.2^{4.15} = 2.1$.
\end{solution}

\begin{solution}{exo:b2:eruptive-styles:4}
$V_{\mathrm{DRE}} = 0.29\times1000/2500 = \qty{0.116}{km^3}$; $m =
\num{0.29e9}\times1000 = \qty{2.9e11}{kg}$; $M = \log_{10}(\num{2.9e11}) - 7 = 11.46
- 7 = 4.46$.
\end{solution}

\begin{solution}{exo:b2:eruptive-styles:5}
$k = \ln(40/10)/(15 - 5) = 1.386/10 = \qty{0.139}{km^{-1}}$; $T_0 = 40\,e^{0.139\times5}
= 40\times2.00 = \qty{80}{cm}$; $\ln2/k = \qty{5.0}{km}$. Avec $T_0 =
\qty{8.0e-4}{km}$: $V = 2\times\num{8.0e-4}/0.139^2 = \qty{0.083}{km^3} =
\qty{8.3e7}{m^3}$, VEI~3.
\end{solution}

\begin{solution}{exo:b2:eruptive-styles:6}
(a) $\int_0^x T_0e^{-ky}2y\,dy = 2T_0\big[-\tfrac{y}{k}e^{-ky} - \tfrac{1}{k^2}e^{-ky}\big]_0^x
= \tfrac{2T_0}{k^2}[1 - (1 + kx)e^{-kx}]$. (b) En itérant $u = \ln(2(1 + u))$ à partir de
$u = 1$~: $1.386$, $1.564$, $1.635$, $1.662$, $1.672$, \dots, convergeant vers
$u = 1.678$, pour lequel $(1 + u)e^{-u} = 2.678\times0.1867 = 0.500$. (c) $x_{1/2} =
1.678/0.139 = \qty{12.1}{km}$, à l'intérieur de l'isopaque le plus externe ($\sqrt A =
\qty{15}{km}$)~; pourtant le volume au-delà de cet isopaque est $(1 + 2.09)e^{-2.09} =
\qty{38}{\percent}$ du total. Près de deux cinquièmes du dépôt se trouvent là où
rien n'a été mesuré~: le volume repose sur l'extrapolation.
\end{solution}

\begin{solution}{exo:b2:eruptive-styles:7}
(a) $p_s = (5.0/0.444)^2 = \qty{127}{MPa}$, à $127/0.025 = \qty{5.1}{km}$.
(b) $x = (5.0 - 0.444\sqrt{10})/100 = (5.0 - 1.404)/100 = 0.0360$;
$V_g = 0.0360\times8.314\times1123/(0.018015\times10^7) = \qty{1.86e-3}{m^3}$;
$V_m = 0.964/2300 = \qty{4.19e-4}{m^3}$; $\phi = 1.86/(1.86 + 0.419) = 0.82$.
(c) Un liquide basaltique est des milliers à des milliards de fois moins visqueux qu'un
liquide rhyolitique (\cref{ch:b2:magma-emplacement})~: ses bulles remontent,
fusionnent et s'échappent, et le gaz quitte le \omterm{def:b2:magma-emplacement:magma}{magma} à travers une mousse perméable ou à l'évent,
si bien que le \omterm{def:b2:magma-emplacement:magma}{magma} se dégaze en système ouvert au lieu de garder son gaz~; sa
fragmentation, lorsqu'elle se produit, reste confinée à une fontaine près de l'évent. Une
remontée plus lente laisse davantage de temps pour cela.
\end{solution}

\begin{solution}{exo:b2:eruptive-styles:8}
(a) $u_z(0) = 0.75\times10^7/(\pi\times(5000)^2) = \qty{0.095}{m} = \qty{95}{mm}$.
(b) $z_0 = 3.8/0.766 = \qty{5.0}{km}$. (c) $(1 + 1)^{-3/2} = 0.35$ du
maximum.
\end{solution}

\begin{solution}{exo:b2:eruptive-styles:9}
La portée maximale est obtenue à \ang{45}~: $R = v^2/g = 100^2/9.81 = \qty{1.0}{km}$, et
$200^2/9.81 = \qty{4.1}{km}$. La résistance de l'air raccourcit le vol~: ce sont donc
des bornes supérieures, proches de la réalité pour les gros blocs et bien trop longues pour les petits.
Puisque $R = v^2/g$, la zone de danger croît comme le carré de la vitesse de lancement~:
doubler la vitesse de l'explosion quadruple son rayon.
\end{solution}

\begin{solution}{exo:b2:eruptive-styles:10}
(a) Avec $T = T_0e^{-k_1x}$ pour $x < X$ et $T = T_Xe^{-k_2(x - X)}$ au-delà, $T_X =
T_0e^{-k_1X}$ (continuité), utiliser $\int_0^X ye^{-ky}dy = \frac{1}{k^2} - e^{-kX}\big(\frac{X}{k} +
\frac{1}{k^2}\big)$ et $\int_X^\infty ye^{-ky}dy = e^{-kX}\big(\frac{X}{k} +
\frac{1}{k^2}\big)$~:
la partie proximale donne $2T_0\big[\frac{1}{k_1^2} - e^{-k_1X}\frac{k_1X + 1}{k_1^2}\big]$ et la
partie distale $2T_Xe^{k_2X}\int_X^\infty ye^{-k_2y}dy = 2T_X\frac{k_2X + 1}{k_2^2}$~; avec $T_X =
T_0e^{-k_1X}$, leur somme est la formule de l'encadré Méthode.
(b) $T_0 = \qty{0.01285}{km}$~: $2T_0/k_1^2 = 0.0257/0.1332 = \qty{0.193}{km^3}$~;
$e^{-0.365\times8.56} = 0.0440$~; $(0.138\times8.56 + 1)/0.138^2 = 114.5$ et
$(0.365\times8.56 + 1)/0.365^2 = 30.96$~; la correction vaut $0.0257\times0.0440\times83.6
= \qty{0.094}{km^3}$ et $V = \qty{0.287}{km^3}$, soit les \qty{0.29}{km^3} publiés.
Au-delà de $X$ se trouve $2T_X(X/k_2 + 1/k_2^2) = 2\times\num{5.65e-4}\times114.5 =
\qty{0.129}{km^3}$, soit \qty{45}{\percent}.
\end{solution}

\begin{solution}{exo:b2:eruptive-styles:11}
(a) $100^{1/4} = 3.16$~; les lois empiriques donnent $100^{0.241} = 3.03$ et
$100^{0.226} = 2.83$. (b) $H \propto N^{-3/4}$~: une atmosphère plus fortement stratifiée
freine la colonne plus bas~; pour $N$ supérieur de \qty{20}{\percent},
$1.2^{-0.75} = 0.87$, soit une hauteur inférieure de \qty{13}{\percent}. (c) Les panaches faibles sont
courbés par le vent, ce qui les abaisse, tandis que l'air humide libère de la chaleur latente en
leur sein, ce qui les élève~; ces effets pèsent différemment sur les petites et les grandes
éruptions et ne figurent pas dans la loi simple. Et les débits ne sont connus qu'à
un facteur de plusieurs près, une erreur sur la variable portée en abscisse, qui biaise
une pente ajustée vers zéro.
\end{solution}

\begin{solution}{exo:b2:eruptive-styles:12}
(a) $132/400\times100 = 33$ par siècle avant 1900~; $95/125\times100 = 76$ par
siècle depuis. (b) $33/76 = \qty{43}{\percent}$. (c) $21/56 = \qty{38}{\percent}$~:
le même sous-enregistrement, d'un facteur deux à trois. (d) Les grandes éruptions se voient
de plus loin et laissent des dépôts plus épais et plus étendus, qui résistent mieux à l'érosion
plus longtemps~; une éruption de VEI~4 sur une île isolée en 1700 pourrait passer inaperçue, une
éruption de VEI~6 difficilement.
\end{solution}

\begin{solution}{pb:b2:eruptive-styles:1}
\textbf{1.} $\ln T = 5.298$, $4.605$, $3.871$, $2.639$, $1.386$. Les pentes entre
isopaques successifs valent $-0.0693$, $-0.0367$, $-0.0308$ et $-0.0313$ par km~: une
pente raide entre 10 et \qty{20}{km}, une pente douce de 40 à
\qty{120}{km}~; la rupture se situe entre 20 et \qty{40}{km}.

\textbf{2.} $k_1 = \ln(200/100)/10 = \qty{0.0693}{km^{-1}}$; $T_0 = 200\,e^{0.693} =
\qty{400}{cm}$.

\textbf{3.} $167e^{-1.244} = \qty{48.1}{cm}$, $167e^{-2.488} = \qty{13.9}{cm}$,
$167e^{-3.732} = \qty{4.00}{cm}$~: à moins de \qty{1}{\percent} de 48, 14 et 4.0.

\textbf{4.} $400e^{-0.0693X} = 167e^{-0.0311X}$ donne $X = \ln(400/167)/(0.0693 -
0.0311) = 0.873/0.0382 = \qty{22.9}{km}$.

\textbf{5.} $\ln2/0.0693 = \qty{10.0}{km}$ et $\ln2/0.0311 = \qty{22.3}{km}$.

\textbf{6.} $T_0 = \qty{0.0040}{km}$: $V = 2\times0.0040/0.0693^2 =
\qty{1.67}{km^3}$.

\textbf{7.} $e^{-0.0693\times22.9} = 0.205$; $(0.0311\times22.9 + 1)/0.0311^2 = 1770$;
$(0.0693\times22.9 + 1)/0.0693^2 = 539$; correction $2\times0.0040\times0.205\times1231 =
\qty{2.02}{km^3}$~; $V = 1.67 + 2.02 = \qty{3.68}{km^3}$, dont la correction représente
\qty{55}{\percent}.

\textbf{8.} $3.68/1.67 = 2.2$. Les cendres distales fines sont minces, facilement érodées,
remaniées ou tombées à la mer~; dans un vieux dépôt, elles sont souvent perdues ou jamais
cartographiées, et le volume est alors sous-estimé d'un facteur deux ou plus.

\textbf{9.} $\qty{3.68e9}{m^3}$ se situe entre $10^9$ et $10^{10}$~: VEI~5.

\textbf{10.} $V_{\mathrm{DRE}} = 3.68\times900/2400 = \qty{1.38}{km^3}$.

\textbf{11.} $m = \num{3.68e9}\times900 = \qty{3.31e12}{kg}$.

\textbf{12.} $M = \log_{10}(\num{3.31e12}) - 7 = 12.52 - 7 = 5.52$.

\textbf{13.} $3.68/1.3 = 2.8$ fois.

\textbf{14.} Une colonne plus haute que \qty{25}{km}, et un style plinien.

\textbf{15.} $Q = \num{3.31e12}/36000 = \qty{9.2e7}{kg/s}$; $I = 7.96 + 3 = 11.0$.

\textbf{16.} $0.304\times(\num{9.2e7})^{0.241} = 0.304\times83.0 = \qty{25.2}{km}$.

\textbf{17.} $0.345\times(\num{9.2e7})^{0.226} = 0.345\times63.1 = \qty{21.8}{km}$.

\textbf{18.} La hauteur mesurée de \qty{31}{km} dépasse de \qty{23}{\percent} la première
prévision et de \qty{42}{\percent} la seconde, ce qui reste dans la
dispersion des lois, de \qtyrange{53}{60}{\percent}~: cohérent.

\textbf{19.} $Q = (31/0.304)^{1/0.241} = 102.0^{4.149} = \qty{2.2e8}{kg/s}$, soit 2.3 fois
le débit moyen. Les fragments enregistrent la colonne la plus haute, au pic de
l'éruption, tandis que la question~15 moyenne sur toute la phase~; une partie de la masse
peut manquer au dépôt (cendres distales, \omterm{def:b2:eruptive-styles:pdc}{écoulements pyroclastiques}).

\textbf{20.} Oui~: de 22 à \qty{31}{km}, bien au-dessus de la tropopause, à \qty{16}{km}.
Le gaz porté dans la stratosphère se répand autour du globe~; le gaz clé est
le dioxyde de soufre, et les vingt millions de tonnes que le Pinatubo y a injectées en 1991
ont refroidi le globe d'environ un demi-degré pendant deux ans.

\textbf{21.} $\lambda_3 = 399/75 = 5.32$, $\lambda_4 = 95/125 = 0.760$, $\lambda_5 = 22/225 =
0.0978$, $\lambda_6 = 8/525 = 0.0152$ par an~; logarithmes $0.726$, $-0.119$, $-1.010$,
$-1.817$.

\textbf{22.} La droite donne $10^{0.73} = 5.37$, $10^{-0.12} = 0.759$, $10^{-0.97} =
0.107$ et $10^{-1.82} = 0.0151$~: à moins de \qty{1}{\percent}, \qty{0}{\percent},
\qty{10}{\percent} et \qty{1}{\percent} des décomptes.

\textbf{23.} Le catalogue des grandes éruptions reste complet plus loin dans le temps,
car on les remarque de plus loin et leurs dépôts survivent plus longtemps~; une fenêtre
commençant en 1500 sous-compterait gravement les éruptions de VEI~3, tandis qu'une fenêtre
commençant en 1950 ne contiendrait presque aucune éruption de VEI~6.

\textbf{24.} $\lambda_5 = \num{0.107}$ par an, une période de retour de $1/0.107 =
\num{9.3}\text{ ans}$~; $1 - e^{-1.07} = 1 - 0.343 = \qty{66}{\percent}$ en dix ans.

\textbf{25.} $\lambda_7 = 10^{3.28 - 5.95} = 10^{-2.67} = 0.00214$ par an, soit une tous les
470 ans~; le catalogue donne $2/1025 = 0.00195$ par an~: accord à moins de
\qty{10}{\percent}, pour un point qui n'a pas servi à l'ajustement.

\textbf{26.} $\lambda_6 = 10^{3.28 - 5.10} = 10^{-1.82} = 0.0151$ par an~; $\lambda t = 1.51$
et $1 - e^{-1.51} = 1 - 0.220 = 0.780$~: la probabilité d'au moins une éruption
de VEI~6 ou plus quelque part sur Terre au cours des cent prochaines années est
\textbf{\qty[text-series-to-math=true]{78}{\percent}}. Elle repose sur un taux constant,
sur des éruptions indépendantes les unes des autres (un processus de Poisson), et sur un
catalogue complet dans ses fenêtres.
\end{solution}
